\textbf{1. Kiến thức cần nhớ}
\begin{itemize}
\item \textbf{Bất đẳng thức} là hệ thức có dạng $a < b$ (hoặc $a > b, a \le b, a \ge b$).
\item \textbf{Tính chất bắc cầu}: Nếu $a < b$ và $b < c$ thì $\boxed{a < c}$.
\item \textbf{Liên hệ giữa thứ tự và phép cộng}: Khi cộng cùng một số vào cả hai vế của một bất đẳng thức, ta được bất đẳng thức mới cùng chiều:
$$\boxed{a < b \implies a + c < b + c}$$
\item \textbf{Liên hệ giữa thứ tự và phép nhân}:
\begin{itemize}
\item Khi nhân cả hai vế với một số \textbf{dương}, bất đẳng thức \textbf{giữ nguyên chiều}:
$$\boxed{a < b\text{ và } c > 0 \implies ac < bc}$$
\item Khi nhân cả hai vế với một số \textbf{âm}, bất đẳng thức \textbf{đổi chiều}:
$$\boxed{a < b\text{ và } c < 0 \implies ac > bc}$$
\end{itemize}
\end{itemize}

\textbf{2. Phương pháp giải các dạng}

%% DANG 9D2?5-1
\textbf{Dạng 1. Nhận biết, so sánh, viết bất đẳng thức}
\begin{enumerate}
\item Dùng kí hiệu $\le$ cho các cụm từ: ``không lớn hơn'', ``nhiều nhất là'', ``tối đa''.
\item Dùng kí hiệu $\ge$ cho các cụm từ: ``không nhỏ hơn'', ``ít nhất là'', ``tối thiểu''.
\item Để so sánh hai số $a$ và $b$, xét hiệu $a - b$:
\begin{itemize}
\item Nếu $a - b > 0 \implies a > b$.
\item Nếu $a - b < 0 \implies a < b$.
\end{itemize}
\end{enumerate}

%% DANG 9D2?5-2
\textbf{Dạng 2. Liên hệ giữa thứ tự với phép cộng, phép nhân}
\begin{enumerate}
\item Vận dụng các quy tắc cộng cùng một số và nhân với cùng một số để biến đổi bất đẳng thức:
$$\text{Từ } a < b \implies ma + n < mb + n\quad (\text{nếu } m > 0)$$
$$\text{Từ } a < b \implies ma + n > mb + n\quad (\text{nếu } m < 0)$$
\item Lưu ý: Nhớ đổi chiều bất đẳng thức khi nhân hoặc chia hai vế cho số âm.
\end{enumerate}

%% DANG 9D2?5-3
\textbf{Dạng 3. Chứng minh bất đẳng thức đơn giản}
\begin{enumerate}
\item \textbf{Phương pháp xét hiệu}: Để chứng minh $A \ge B$, ta xét hiệu $A - B$ và biến đổi thành bình phương của một tổng/hiệu hoặc tổng các bình phương:
$$A - B = (X \pm Y)^2 \ge 0 \implies A \ge B.$$
\item \textbf{Phương pháp biến đổi tương đương}: Biến đổi bất đẳng thức cần chứng minh về một bất đẳng thức luôn đúng (như $X^2 \ge 0$).
\item Luôn chỉ rõ điều kiện để dấu ``='' xảy ra.
\end{enumerate}
